%-------------------------
% Resume in Latex
% Author : Aras Gungore
% License : MIT
%------------------------

\documentclass[letterpaper,11pt]{article}

\usepackage{latexsym}
\usepackage[empty]{fullpage}
\usepackage{titlesec}
\usepackage{marvosym}
\usepackage[usenames,dvipsnames]{color}
\usepackage{verbatim}
\usepackage{enumitem}
\usepackage[hidelinks]{hyperref}
\usepackage{fancyhdr}
\usepackage[english]{babel}
\usepackage{tabularx}
\usepackage{hyphenat}
\usepackage{fontawesome}
\usepackage{graphicx}
\usepackage[T1]{fontenc}
\usepackage[utf8]{inputenc}
\usepackage{enumitem}
\usepackage{lmodern} % Loads a font that supports PDF searching/copying better
\input{glyphtounicode}

%---------- FONT OPTIONS ----------
% sans-serif
% \usepackage[sfdefault]{FiraSans}
% \usepackage[sfdefault]{roboto}
% \usepackage[sfdefault]{noto-sans}
% \usepackage[default]{sourcesanspro}

% serif
% \usepackage{CormorantGaramond}
% \usepackage{charter}

\pagestyle{fancy}
\fancyhf{} % clear all header and footer fields
\fancyfoot{}
\setlength{\footskip}{5pt}
\renewcommand{\headrulewidth}{0pt}
\renewcommand{\footrulewidth}{0pt}

% Adjust margins
\addtolength{\oddsidemargin}{-0.6in}
\addtolength{\evensidemargin}{-0.5in}
\addtolength{\textwidth}{1.2in}
\addtolength{\topmargin}{-.7in}
\addtolength{\textheight}{1.5in}

\urlstyle{same}

\raggedbottom
\raggedright
\setlength{\tabcolsep}{0in}

% Sections formatting
\titleformat{\section}{
  \vspace{-4pt}\scshape\raggedright\large
}{}{0em}{}[\color{black}\titlerule \vspace{-5pt}]

% Ensure that generate pdf is machine readable/ATS parsable
\pdfgentounicode=1

\newcommand{\resumeSubHeadingListStart}{\begin{itemize}[leftmargin=0.15in, label={}]}
\newcommand{\resumeSubHeadingListEnd}{\end{itemize}}
\newcommand{\resumeItemListStart}{\begin{itemize}}
\newcommand{\resumeItemListEnd}{\end{itemize}\vspace{-5pt}}

%-------------------------------------------
%%%%%%  RESUME STARTS HERE  %%%%%%%%%%%%%%%%%%%%%%%%%%%%
\setlist[itemize]{
    itemsep=0.35pt,      % Space between bullets
    topsep=2pt,       % Space between the text and the first bullet
    parsep=2pt,       % Space between paragraphs inside a bullet
    partopsep=0pt
}
\begin{document}

%---------- HEADING ----------
\begin{minipage}[c]{0.7\textwidth}
    {\Huge \scshape Ashwin Balaji Arun Kumar} \\ \vspace{1pt}

    \small
    \textbf{E-mail:} \hspace{7pt} \href{mailto: arunkumar.ashwin@outlook.com}{arunkumar.ashwin@outlook.com} \hspace{-5pt}
    \quad 
    \textbf{Telefon:} \href{tel:+4917665908289}{+49 17665908289} \\
    \textbf{LinkedIn:} \href{https://www.linkedin.com/in/ashwinbalaji-ar/}{linkedin.com/in/ashwinbalaji-ar}
    \quad 
    \textbf{Github:} \href{https://github.com/ashwin-ar/}{github.com/ashwin-ar/} \\
    \quad 
    \textbf{Standort:} Dortmund, Germany
\end{minipage}
\hspace{0.07\textwidth} % This creates your LEFT indent
\begin{minipage}[c]{0.2\textwidth}
    \flushright
    \vspace{0pt} 
    % These two lines control the border thickness and gap
    \setlength{\fboxsep}{0pt}  % No gap between photo and border
    \setlength{\fboxrule}{0.5pt} % Thickness of the line
    \fbox{\includegraphics[trim=14 9 14 0, clip, width=2.4cm, keepaspectratio]{{Ashwin_passport}.jpg}}
\end{minipage}

%----------- EDUCATION -----------
\vspace{-5pt}
\section{Education}
\noindent
\textbf{Fachhochschule Dortmund} \hfill {\small Dortmund, DE} \\
{\small \textit{Master of Engineering (M.Eng.) in Eingebettete Systeme Technik} \hfill \small {10/2025 \textbf{--} Present} \\ \vspace{-1pt}
\textit{(Ergebnisse des ersten Semesters stehen noch aus)}}
\vspace{3pt}

\noindent
\textbf{Anna University} \hfill {\small Coimbatore, Indien} \\
{\small \textit{Bachelor of Engineering (B.E.) in Elektrotechnik und Elektronik} \hfill \small {08/2018 \textbf{--} 06/2022} \\
\textit{GPA: 1.5 / 4.0 (Deutsches Äquivalent)}}

%----------- SKILLS -----------
\vspace{-3pt}
\section{Skills}
\vspace{2pt}

  \textbf{Sprachen \& Tools:} {C, C++, Python, Assembly(Grundlagen), Matlab/Simulink, Git, SVN} \\ \vspace{2pt}
  \textbf{Hardware:} {TI C2000 (TMS320F), Renesas RL78, STM32 (F4 Series), ESP32, RP2040} \\ \vspace{2pt}
  \textbf{Kommunikationsprotokolle:} {I2C, UART, SPI, Modbus, RS485, MQTT} \\ \vspace{2pt}
  \textbf{Leistungselektronik:} {Frequenzumrichter, MPPT, AC/DC-Wandler, PV-Anlagen, Batteriemanagements} \\ \vspace{2pt} 
  \textbf{Laborgeräte:} {Multimeter, Hochfrequenz-Digital-Oszilloskop, Frequenzgenerator, PV-Simulatoren} \\ \vspace{2pt}
  %\textbf{Qualität \& Standards:} {Statische Analyse(Coverity), Fehlermöglichkeits und Einflussanalyse, IEC-Normen} \\ \vspace{2pt}
  %\textbf{Architektur \& KI:} {Microsoft 365 (Excel-Experte), KI-gestützte Entwicklung (LLM-APIs, Gemini, OpenAI)} \\ \vspace{2pt}
  %\textbf{Build Systeme:} {Latex, Linker Skripte, Make, CMake, Makefiles} \\ \vspace{2pt}

% --- Work Experience ---
\vspace{-8pt}
\section{Professional Experience}
\textbf{Schneider Electric} \hfill Mumbai, Indien \\
\textit{Stellvertretender Leiter \textbf{--} Firmware} \hfill 08/2023  \textbf{--}  09/2025 \\
\begin{itemize}[leftmargin=*]
    \item Leitung der Firmware Entwicklung für \textbf{3–20 HP Frequenzumrichter} auf TI \& RL78-MCUs.
    \item Entwurf eine neue modularen \textbf{Applikationslayer mit HAL} zur Entkopplung von Steuerung und Hardware.
    \item Implementierung eines hocheffizienten (98 \%) \textbf{P\&O-MPPT} Algorithmus für Gleichstrom-Solarpumpensysteme.   
    %\item Erfahrung in der Steuerung von \textbf{3-Phasen- und Einphasen-Motoren} (Induktion, BLDC) sowie PID-Regel.
    \item Optimierung des \textbf{Kurzschlussschutzes} auf \textbf{200 ns} via Trip-Zones zur \textbf{Reduzierung von Feldausfällen}.
    \item \textbf{CI/CD-Integration statischer Analysen} zur Gewährleistung von 10+ hochwertigen Firmware Releases.
    \item \textbf{EEPROM-Emulation auf Flash-Basis} als Hardware-Ersatz zur deutlichen Senkung der Stückkosten.    
    %\item \textbf{Betreuung einen Hochschulabsolventen} im Entwicklung von TI treibern und Git workflows.
    %\item Entwicklung eines \textbf{Python-UART Telemetrie Tool} zur Behebung von Leistungsproblemen.
    %\item \textbf{Startfehler behoben} durch \textbf{I2C-Bus-Recovery} bei Arbitrierungs-Hangs
    %\item \textbf{Modellierte} Antriebsregelkreise in \textbf{Matlab/Simulink} zur Überprüfung der PID-Stabilität.  
  \end{itemize}
\vspace{8pt}
\noindent
\textit{Graduate Engineer Trainee \textbf{--} Electromechanical} \hfill 07/2022 \textbf{--} 08/2023 \\
\begin{itemize}[leftmargin=*]
    \item Produktpflege der \textbf{6A -- 100A RCBO-Reihe} durch Designänderungen und Reklamationsmanagement.  
    \item Reduziert der \textbf{RCBO-Kurzschlussenergie um 4 \%} (gem. IEC 61009-1) durch Kontaktoptimierung.
    %\item \textbf{Temperatursenkung um 7°C} bei Lasttrennschaltern (gem. IEC 60947-3) durch thermische Optimierung.
    \item Identifizierung von über 30 Fehlerfällen in der DFMEA und \textbf{Behebung von 2 Fehlern} mit RPN>144.
  \end{itemize}
\vspace{8pt}
\noindent
\textbf{Virtual Sapiens, Coimbatore, Indien} \textit{- Hardware Intern} \hfill 08/2020 \textbf{--} 01/2022
\begin{itemize}[leftmargin=*]
\item Bau eines Telepresence-Bots mit \textbf{Raspberry Pi für Low-Latency-Videostreaming (<800ms)} via HTTP und \textbf{AWS MQTT} zur Bot-Steuerung.
\item Entwicklung eines Regelalgorithmus für BLDC-Motoren auf Basis des STM32F4103.
\item Implementierung der \textbf{I2C-Kommunikation} zwischen Raspberry Pi und STM32 zum Lenkbefehlsaustausch.
\end{itemize}

% --- Projects ---
\section{Projects}
\textbf{i-Practace \textbf{--} Automatic Tennis ball throwing machine}
\begin{itemize}[leftmargin=*]
  %\item Gefördert durch \textbf{Indischen Ministerium für Wiss. \& Tech.} mit einem Betrag von \textbf{9.900 Euro} (9,48L INR).
  \item PCB-Design mit Ki-Cad zur Integration eines Dimmerschaltkreises und eines SMPS.
  \item Hardwareauswahl für Motoren/Sensoren sowie Leitung der gesamten POC-Entwicklung bis zur Feldtestreife.
\end{itemize}

% --- Patents & Publications ---
% \section{Patents and Publications}
% \textbf{}
% \begin{itemize}[leftmargin=*]
%     \item \textbf{Automatisiertes Tennis Spieltrainingsgerät} - \small Indisches Geschmacksmuster Nr. - 367774-001 (Nr. 20/2023)
%     \item \href{https://ieeexplore.ieee.org/document/9675583} {Vergleichende Analyse industrieller IoT Kommunikationsprotokolle und ihrer zukünftigen Ausrichtung} - IEEE XPLORE
% \end{itemize}
%----------- Languages -----------

\section{Languages}
  \vspace{2pt}
  \textbf{German:} {A2 (wird aktiv verbessert)} 
  \quad
  \textbf{English:} {C1}


\end{document}